% HRAI-GO Platform | Documento pedagogico de arquitetura e governanca
% Fonte editavel. Compilacao: xelatex (duas passagens).
\documentclass[11pt,a4paper,oneside]{article}
\usepackage[brazil]{babel}
\usepackage{fontspec}
\setmainfont{Noto Serif}
\setsansfont{Noto Sans}
\setmonofont[Scale=.83]{DejaVu Sans Mono}
\usepackage{geometry}
\geometry{left=24mm,right=23mm,top=24mm,bottom=23mm,headheight=15pt}
\usepackage{microtype}
\usepackage{setspace}
\setstretch{1.13}
\usepackage{enumitem}
\setlist[itemize]{leftmargin=16pt,itemsep=3pt,topsep=4pt}
\setlist[enumerate]{leftmargin=18pt,itemsep=3pt,topsep=4pt}
\usepackage[table]{xcolor}
\definecolor{ink}{HTML}{152739}
\definecolor{muted}{HTML}{526579}
\definecolor{blue}{HTML}{1664A4}
\definecolor{teal}{HTML}{00877B}
\definecolor{pale}{HTML}{EAF3F7}
\definecolor{light}{HTML}{F4F7FA}
\definecolor{amber}{HTML}{9F6200}
\usepackage{titlesec}
\titleformat{\section}{\Large\bfseries\sffamily\color{ink}}{\thesection}{.8em}{}
\titleformat{\subsection}{\large\bfseries\sffamily\color{blue}}{\thesubsection}{.7em}{}
\titleformat{\subsubsection}{\normalsize\bfseries\sffamily\color{ink}}{\thesubsubsection}{.7em}{}
\titlespacing*{\section}{0pt}{17pt}{9pt}
\titlespacing*{\subsection}{0pt}{12pt}{5pt}
\usepackage{fancyhdr}
\pagestyle{fancy}
\fancyhf{}
\fancyhead[L]{\sffamily\footnotesize\color{muted} HRAI-GO | Plataforma de Inteligência Regulatória}
\fancyhead[R]{\sffamily\footnotesize\color{muted} Documento pedagógico}
\fancyfoot[L]{\sffamily\footnotesize\color{muted} Proposta conceitual \textbullet\ 11 nov. 2026}
\fancyfoot[R]{\sffamily\footnotesize\color{muted}\thepage}
\renewcommand{\headrulewidth}{0.35pt}
\renewcommand{\footrulewidth}{0pt}
\usepackage{tabularx,longtable,booktabs,array,multirow}
\renewcommand{\arraystretch}{1.22}
\newcolumntype{Y}{>{\raggedright\arraybackslash}X}
\usepackage[most]{tcolorbox}
\newtcolorbox{keybox}[1][]{colback=pale,colframe=blue!55!white,boxrule=.7pt,arc=2mm,left=11pt,right=11pt,top=9pt,bottom=9pt,#1}
\newtcolorbox{notebox}[1][]{colback=light,colframe=teal!65!white,boxrule=.55pt,arc=2mm,left=10pt,right=10pt,top=8pt,bottom=8pt,#1}
\newtcolorbox{warnbox}[1][]{colback=orange!5,colframe=amber!65,boxrule=.65pt,arc=2mm,left=10pt,right=10pt,top=8pt,bottom=8pt,#1}
\usepackage{tikz}
\usetikzlibrary{arrows.meta,positioning,calc,fit,shapes.geometric,matrix}
\usepackage{float}
\usepackage{caption}
\captionsetup{font=small,labelfont=bf}
\usepackage{listings}
\usepackage{seqsplit}
\usepackage{graphicx}
\lstdefinestyle{code}{basicstyle=\ttfamily\footnotesize,columns=fullflexible,breaklines=true,breakatwhitespace=false,keepspaces=true,showstringspaces=false,frame=single,framesep=7pt,rulecolor=\color{blue!30},backgroundcolor=\color{light},aboveskip=8pt,belowskip=8pt,upquote=true}
\lstset{style=code}
\usepackage{hyperref}
\hypersetup{unicode=true,colorlinks=true,linkcolor=blue,urlcolor=teal,citecolor=blue,pdfauthor={HRAI-GO},pdftitle={HRAI-GO: Human-centred Regulatory AI, Governance Observatory},pdfsubject={Plataforma de inteligencia regulatoria humano-IA e Shadow Council}}
\usepackage{bookmark}
\usepackage{csquotes}
\usepackage{amssymb,amsmath}
\usepackage{needspace}
\newcommand{\term}[1]{{\ttfamily\small\seqsplit{#1}}}
\newcommand{\eng}[1]{\textit{#1}}
\newcommand{\nsource}{\textbf{Fonte:} }
\newcommand{\nproposal}{\textbf{Proposta:} }
\newcommand{\nexample}{\textbf{Exemplo fictício:} }
\newcommand{\tinytag}[1]{{\sffamily\footnotesize\color{teal}\textbf{#1}}}
\begin{document}
\begin{titlepage}
\thispagestyle{empty}
\vspace*{12mm}
{\sffamily\bfseries\color{teal}\Large HRAI-GO / CONCEPT SERIES\par}
\vspace{2mm}
{\sffamily\color{muted}\large Human-centred Regulatory AI, Governance Observatory\par}
\vspace{12mm}
{\sffamily\bfseries\color{ink}\fontsize{29}{34}\selectfont Plataforma Aberta de\\Inteligência Regulatória\\e Governança de IA\par}
\vspace{9mm}
{\sffamily\fontsize{18}{23}\selectfont\color{blue} Ontologias, agentes contínuos, avaliação de conformidade, comparação internacional e \textit{Shadow Council}\par}
\vspace{16mm}
\begin{tikzpicture}[x=1cm,y=1cm,>=Latex]
\fill[blue!7,rounded corners=7pt] (0,0) rectangle (15.4,5.9);
\node[draw=blue!70,fill=white,rounded corners=5pt,minimum width=5.3cm,minimum height=.95cm,font=\sffamily\bfseries\large,text=ink] (kg) at (7.7,.85) {Canonical Knowledge Graph};
\node[draw=teal!70,fill=white,rounded corners=4pt,minimum width=6.6cm,minimum height=.9cm,font=\sffamily\bfseries,text=ink] (agents) at (7.7,4.9) {Multi-Agent Orchestration};
\node[draw=blue!60,fill=white,rounded corners=3pt,minimum width=3.3cm,minimum height=1.15cm,align=center,font=\sffamily\small\bfseries] (c) at (2.55,3.05) {Compliance\\Intelligence};
\node[draw=blue!60,fill=white,rounded corners=3pt,minimum width=3.3cm,minimum height=1.15cm,align=center,font=\sffamily\small\bfseries] (p) at (7.7,3.05) {Policy\\Intelligence};
\node[draw=blue!60,fill=white,rounded corners=3pt,minimum width=3.3cm,minimum height=1.15cm,align=center,font=\sffamily\small\bfseries] (s) at (12.85,3.05) {Shadow Council\\Deliberation};
\draw[->,very thick,teal!75] (agents.south) -- (p.north);
\draw[->,thick,blue!75] (kg.north -| c) -- (c.south);
\draw[->,thick,blue!75] (kg.north) -- (p.south);
\draw[->,thick,blue!75] (kg.north -| s) -- (s.south);
\draw[<->,thin,dashed,gray] (c.east) -- (p.west);
\draw[<->,thin,dashed,gray] (p.east) -- (s.west);
\end{tikzpicture}
\vfill
{\sffamily\bfseries\color{ink} Documento pedagógico de consolidação e especificação conceitual\par}
\vspace{3mm}
{\sffamily\color{muted} Versão 0.3 | 17 de novembro de 2026 | Português brasileiro\par}
\vspace{4mm}
\begin{minipage}{.94\textwidth}
\small\color{muted} Conteúdo: síntese dos três estudos fornecidos, arquitetura proposta e funcionalidade nova \textit{Shadow Council}. Exemplos fictícios, regras ilustrativas e propostas de engenharia são identificados como tais. Este material não substitui interpretação jurídica oficial nem comprova representatividade social.
\end{minipage}
\end{titlepage}

\pagenumbering{roman}
\tableofcontents
\clearpage
\pagenumbering{arabic}

\section{Resumo executivo: o que estamos construindo?}
\begin{keybox}
\textbf{Ideia central.} O HRAI-GO deixa de ser apenas uma ontologia de sistemas de IA de alto risco e evolui para uma \textbf{infraestrutura aberta, multinacional e multiagente de inteligência regulatória}. Ela integra duas tarefas diferentes: avaliar uma aplicação frente às normas aplicáveis e avaliar criticamente a própria adequação da regulação às necessidades de uma sociedade. Uma terceira capacidade, o \textit{Shadow Council}, experimenta deliberações representativas de perspectivas documentadas dos diferentes stakeholders, sem fingir que agentes sintéticos são pessoas reais ou mandatos democráticos.
\end{keybox}

\textbf{Nome.} HRAI-GO significa \textit{Human-centred Regulatory AI, Governance Observatory}: uma plataforma de inteligência regulatória humano--IA. A sigla nasceu como \textit{High-Risk AI}, sentido que permanece apenas nos namespaces técnicos das ontologias (\texttt{hrai-*.owl}).

\subsection{As quatro capacidades da plataforma}
\begin{center}
\begin{tabularx}{\textwidth}{@{}p{.245\textwidth}Y@{}}
\toprule
\textbf{Componente} & \textbf{Pergunta que responde}\\\midrule
\term{NORM} (Normative Intelligence) & O que foi proposto, aprovado, entrou em vigor, é aplicável, foi alterado ou revogado em cada jurisdição?\\
\term{VERIFY} (Application Compliance) & Quais obrigações alcançam uma aplicação específica, que evidências existem e que questões ainda não foram resolvidas?\\
\term{DEBATE} (Policy Intelligence + Shadow Council) & Que tensões, omissões, argumentos e perspectivas devem entrar no debate público sobre a qualidade da regulação?\\
\term{FORESIGHT} (Policy Alternatives) & Que cenários e efeitos condicionais das alternativas regulatórias podem ser investigados e testados?\\
\bottomrule
\end{tabularx}
\end{center}

\subsection{Os resultados esperados}
\begin{itemize}
\item Uma \textbf{biblioteca comparativa dinâmica} entre legislações, políticas e modelos regulatórios de países diferentes, incluindo equivalência conceitual e diferenças jurídicas.
\item Uma \textbf{ontologia OWL 2} canônica, instanciada continuamente por fatos com proveniência e temporalidade.
\item Um \textbf{motor de compliance} que combine raciocínio ontológico, regras normativas e validação de evidências.
\item Uma \textbf{rede aberta de agentes} de reguladores, especialistas, empresas, sociedade civil e grupos diretamente afetados.
\item Um \textbf{Shadow Council} que represente a diversidade de posições documentadas e teste estruturas de debate, incluindo o registro explícito de dissenso.
\item Um \textbf{observatório de tensões e lacunas} regulatórias, com investigações reprodutíveis e propostas contextualizadas.
\end{itemize}

\begin{warnbox}
\textbf{Separação epistemológica obrigatória.} (1) Texto jurídico oficial, (2) interpretação jurídica, (3) evidência empírica, (4) hipótese de agente, (5) posição de stakeholder e (6) proposta normativa são objetos diferentes. Nenhum dos seis deve ser silenciosamente convertido em outro.
\end{warnbox}

\section{Origem das ideias, limites das fontes e método de leitura}
\subsection{O que vem dos materiais fornecidos}
A proposta toma como base três documentos anexados pelo solicitante:\par
\begin{enumerate}
\item \eng{Guia de Conformidade e Governança para Sistemas de Inteligência Artificial de Alto Risco (HRAI)}: classificação de risco, governança de dados, supervisão humana, auditoria, evidências, mecanismos de controle e monitoramento contínuo \cite{guia}.
\item \eng{Análise Comparativa e Diretrizes Estratégicas para a Governança e Plano de Inteligência Artificial no Brasil}: dilemas de importação de modelos regulatórios, capacidades nacionais, soberania, inclusão, custos e orçamento de governança \cite{comparativo}.
\item \eng{Brainstorming: tipos de artefatos técnicos que não fiquem datados}: matriz dinâmica de riscos, padrões de interoperabilidade, repositório colaborativo de vieses, compras públicas e telemetria de soberania \cite{brainstorm}.
\end{enumerate}
As fontes oferecem o \textbf{problema, os temas e hipóteses de análise}, não uma ontologia já especificada, uma amostra representativa real, evidência causal de que determinada política funciona ou uma validação da plataforma. Os nomes canônicos, mecanismos algorítmicos, regras e protocolos seguintes são \textbf{propostas de modelagem}.

\subsection{O que os estudos não demonstram}
\begin{itemize}
\item As caracterizações gerais de países são pontos de partida analíticos. Não substituem exame de normas, setores, entes subnacionais e períodos de vigência.
\item A proporção de recursos de um eixo orçamentário não demonstra, por si, falta de capacidade fiscalizatória. Exige comparação com gastos executados, pessoal, competências e demanda.
\item A documentação fornecida não permite concluir que perspectivas geradas por IA reproduzem fielmente a distribuição de opiniões dos stakeholders.
\item Afirmações como ``um método de alinhamento resolve explicabilidade jurídica'' devem ser tratadas como hipóteses contestáveis, não como propriedades garantidas da tecnologia.
\end{itemize}

\subsection{Cuidado com o tempo jurídico}
A legislação deve ser consultada em fonte oficial e sob uma data explícita. O \textbf{PL 2338/2023}, por exemplo, é tratado neste documento como \term{LegislativeProposal}, não como lei em vigor: a ficha da Câmara consultada para esta versão o registra em tramitação, aguardando parecer na comissão especial \cite{camara}. Para o EU AI Act, diferentes regras possuem datas distintas de aplicação, como explicita o calendário oficial europeu \cite{eutimeline}. Essas observações são um \textbf{retrato documentado em 11/11/2026}, não um estado imutável.

\section{Um modelo pedagógico: três tipos de conhecimento}
Imagine que a plataforma avalie um sistema fictício de triagem automatizada de currículos no Brasil.
\begin{enumerate}
\item \textbf{Conhecimento jurídico:} a LGPD impõe regras específicas para tratamento de dados pessoais, sob condições e exceções que precisam ser formalmente identificadas. Um projeto de lei sobre IA, por outro lado, é uma proposta e não uma obrigação vigente.
\item \textbf{Conhecimento técnico:} o produto usa dados pessoais, produz ranqueamento e tem ou não mecanismos documentados de revisão humana, testes de vieses e registro de decisões.
\item \textbf{Conhecimento político-regulatório:} pode existir tensão entre custos antecipados de compliance e acessibilidade de pequenas empresas, sem que isso dispense as obrigações efetivamente vigentes.
\end{enumerate}
A plataforma conecta esses fatos, mas não os confunde. A mesma avaliação de conformidade pode ser \textbf{juridicamente válida} sob determinada norma e, ainda assim, originar uma discussão pública sobre se a norma deveria ser aperfeiçoada.

\section{Arquitetura geral e fluxos contínuos}
\begin{figure}[H]
\centering
\resizebox{\linewidth}{!}{%
\begin{tikzpicture}[node distance=6mm and 7mm,>=Latex, every node/.style={font=\sffamily\footnotesize},box/.style={draw=blue!65,rounded corners=3pt,fill=light,minimum height=9mm,align=center,inner sep=5pt},core/.style={draw=teal!80,rounded corners=3pt,fill=teal!9,minimum height=9mm,align=center,inner sep=6pt}]
\node[box,text width=12.0cm] (sources) {Official legal sources \quad | \quad Economic and social indicators \quad | \quad Research \quad | \quad Audits and incidents};
\node[core,text width=12.0cm,below=of sources] (ingest) {Versioned source ingestion, extraction, verification, provenance and review};
\node[core,text width=12.0cm,below=of ingest] (kg) {Canonical Ontologies + Temporal Regulatory Knowledge Graph + National Profiles};
\node[box,text width=3.42cm,below left=8mm and -.3cm of kg.south] (verify) {VERIFY\\Application compliance};
\node[box,text width=3.42cm,right=5mm of verify] (debate) {DEBATE\\Gaps and Shadow Council};
\node[box,text width=3.42cm,right=5mm of debate] (forecast) {FORESIGHT\\Alternatives and scenarios};
\node[core,text width=12.0cm,below=15mm of debate] (users) {Open APIs \quad | \quad Regulatory dashboards \quad | \quad Stakeholder contributions \quad | \quad Audit trails};
\draw[->,thick,teal] (sources)--(ingest);
\draw[->,thick,teal] (ingest)--(kg);
\foreach \n in {verify,debate,forecast} {\draw[->,thick,blue] (kg.south) -- (\n.north); \draw[->,thick,blue] (\n.south)--(users.north);}
\draw[->,dashed,gray] (users.west) to[out=195,in=185,looseness=1.17] (ingest.west);
\end{tikzpicture}%
}
\caption{Arquitetura lógica proposta. As setas representam dependências de informação, não integrações já implementadas.}
\end{figure}

\subsection{Seis camadas operacionais}
\begin{longtable}{@{}p{.29\textwidth}p{.65\textwidth}@{}}
\toprule\textbf{Camada} & \textbf{Função}\\\midrule\endhead
\term{Source Intelligence} & Coleta periódica e por evento, versionamento de páginas, documentos e metadados oficiais.\\
\term{Knowledge Engineering} & Extrai candidatos a fatos, alinha termos, valida semântica e publica instâncias aceitas.\\
\term{Reasoning and Validation} & OWL 2, SWRL, SHACL, SPARQL e regras legais com exceções e hierarquia.\\
\term{Analytical Intelligence} & Avaliações de aplicações, comparações internacionais, indicadores contextuais e lacunas.\\
\term{Deliberative Intelligence} & Agentes de stakeholders, Shadow Council, revisão de argumentos, contradições e dissenso.\\
\term{Open Governance} & APIs, registros de identidade, permissões, auditoria, revisão humana e publicação pública diferenciada.\\
\bottomrule
\end{longtable}

\section{Modelo de conhecimento canônico: OWL, dados e versões}
\subsection{Separar TBox e ABox}
\textbf{TBox} descreve conceitos, axiomas e propriedades: o significado de \term{LegalProvision}, \term{AISystem} ou \term{StakeholderPosition}. \textbf{ABox} contém indivíduos: uma disposição legal identificada, uma instituição, uma avaliação e suas evidências. A publicação de nova norma normalmente atualiza a ABox. Uma nova classe conceitual passa por revisão e versionamento da TBox.

\subsection{Módulos canônicos propostos}
\begin{longtable}{@{}p{.39\textwidth}p{.55\textwidth}@{}}
\toprule\textbf{Arquivo} & \textbf{Cobertura}\\\midrule\endhead
\term{hrai-core.owl} & Entidades nucleares, organizações, aplicações, jurisdições, datas e identidade.\\
\term{hrai-regulatory.owl} & Instrumentos, dispositivos, status legislativo, exceções, normas e obrigações.\\
\term{hrai-risk.owl} & Categorias de risco, impactos, danos, severidade e mitigação.\\
\term{hrai-governance.owl} & Agentes responsáveis, controles, papéis e supervisão humana.\\
\term{hrai-assurance.owl} & Evidências, exames, testes, avaliações e conclusões.\\
\term{hrai-monitoring.owl} & Eventos, incidentes, alterações, drift e ações corretivas.\\
\term{hrai-national-context.owl} & Indicadores e capacidades socioeconômicas e institucionais.\\
\term{hrai-comparative.owl} & Equivalência e diferenças entre jurisdições e instrumentos.\\
\term{hrai-argumentation.owl} & Teses, objeções, suporte documental e tensões.\\
\term{hrai-deliberation.owl} & Conselhos, assentos, rodadas, contribuições e pareceres.\\
\term{hrai-provenance.owl} & Agentes de criação, fontes, revisões e linhagem de evidências.\\
\bottomrule
\end{longtable}

\subsection{Principais classes e propriedades novas}
\begin{tabularx}{\textwidth}{@{}p{.37\textwidth}Y@{}}
\toprule\textbf{Classe} & \textbf{Interpretação}\\\midrule
\term{Jurisdiction} & Ordenamento territorial ou supranacional.\\
\term{LegalInstrument} / \term{LegislativeProposal} & Norma jurídica / proposição legislativa (distintas).\\
\term{LegalProvision} / \term{NormativeRequirement} & Dispositivo textual / requisito interpretado e situado.\\
\term{RegulatoryTension} / \term{RegulatoryGap} & Possível conflito entre objetivos / cobertura ou capacidade insuficiente investigada.\\
\term{PolicyClaim} / \term{CounterArgument} & Tese fundamentada / objeção ou contestação.\\
\term{Stakeholder} / \term{StakeholderPosition} & Ator ou grupo / posição identificável e atribuída.\\
\term{Deliberation} / \term{ShadowCouncilSession} & Processo deliberativo / sessão de agentes sintéticos.\\
\term{SamplingFrame} / \term{StakeholderStratum} & Universo documental ou populacional / estrato usado na composição.\\
\term{CouncilSeat} / \term{PerspectiveAgent} & Assento de uma perspectiva / agente que a simula sob regras explícitas.\\
\term{EvidenceRecord} / \term{Assessment} & Registro verificável / avaliação reprodutível.\\
\bottomrule
\end{tabularx}

\subsection{Relações centrais em inglês}
\begin{lstlisting}
applicableIn(juridicalObject, jurisdiction)
expressesRequirement(legalProvision, normativeRequirement)
subjectToRequirement(aiSystem, normativeRequirement)
substantiatedBy(policyClaim, evidenceRecord)
challenges(counterArgument, policyClaim)
affectsPopulation(policyInstrument, populationGroup)
hasNationalIndicator(jurisdiction, contextIndicator)
representsPerspective(councilSeat, stakeholderStratum)
generatedBy(contribution, perspectiveAgent)
reviewedBy(assessment, reviewer)
usesLegalSnapshot(assessment, legalSnapshot)
\end{lstlisting}
\textbf{Cuidado semântico:} \term{representsPerspective} descreve simulação de uma posição documentada; não significa \term{legallyRepresents}. Representação institucional verdadeira exige autorização verificável.

\subsection{Pequena formalização OWL em Turtle}
\begin{lstlisting}
@prefix hrai: <https://example.org/hrai-go#> .
@prefix owl:  <http://www.w3.org/2002/07/owl#> .
@prefix rdfs: <http://www.w3.org/2000/01/rdf-schema#> .

hrai:LegalInstrument a owl:Class .
hrai:LegislativeProposal a owl:Class .
hrai:NormativeRequirement a owl:Class .
hrai:AISystem a owl:Class .
hrai:StakeholderStratum a owl:Class .
hrai:PerspectiveAgent a owl:Class .
hrai:ShadowCouncilSession a owl:Class .

hrai:subjectToRequirement a owl:ObjectProperty ;
  rdfs:domain hrai:AISystem ;
  rdfs:range hrai:NormativeRequirement .

hrai:representsPerspective a owl:ObjectProperty ;
  rdfs:domain hrai:PerspectiveAgent ;
  rdfs:range hrai:StakeholderStratum .
\end{lstlisting}
\noindent\textit{Nota:} os URIs com \term{example.org} são identificadores ilustrativos. Na implementação, usar domínios institucionais persistentes e políticas explícitas de versionamento de vocabulário.

\section{Atualização contínua das normas e avaliações}
\subsection{Pipeline normativo}
\begin{enumerate}
\item \textbf{Monitoramento:} fontes oficiais por jurisdição, órgão, assunto e instrumento.
\item \textbf{Detecção:} nova publicação, alteração textual, tramitação, eficácia, suspensão, revogação ou retificação.
\item \textbf{Extração candidata:} requisitos, titulares, destinatários, fatos condicionantes, exceções, sanções e prazos.
\item \textbf{Verificação:} ancoragem em trecho original, validação documental e política de revisão proporcional ao risco.
\item \textbf{Publicação:} nova instância/versionamento sem apagar versões anteriores.
\item \textbf{Propagação:} reavaliar só aplicações e análises dependentes de fatos que mudaram.
\end{enumerate}

\subsection{Temporalidade e estado normativo}
Cada objeto regulatório requer, no mínimo, a identificação do documento-fonte, da jurisdição, da versão, do status jurídico e dos intervalos de aplicabilidade. Modelar o tempo sob duas perspectivas:\par
\begin{description}[leftmargin=0pt,style=nextline]
\item[\term{validTime}] Quando o fato ou requisito valia juridicamente para a aplicação.
\item[\term{recordTime}] Quando o sistema conheceu, registrou ou corrigiu aquele fato.
\end{description}
Assim, a pergunta ``o que se sabia em uma auditoria de junho?'' difere de ``quais regras eram aplicáveis em junho?''.

\subsection{Estados controlados}
\begin{lstlisting}
Proposed | UnderDeliberation | Approved |
Enacted | InForce | Applicable | Suspended |
Repealed | Superseded | Uncertain
\end{lstlisting}
\begin{warnbox}
\textbf{Não misturar:} \term{InForce} e \term{Applicable} não são sinônimos. Uma norma pode estar em vigor e conter obrigações com aplicação diferida. Uma proposição legislativa aprovada em apenas uma casa não equivale a uma lei promulgada.
\end{warnbox}

\section{Motor VERIFY: avaliação contínua de aplicações}
\subsection{O dossiê técnico de uma aplicação}
Um \term{ApplicationDossier} documentaria finalidade, contexto, operadores, destinatários, fluxos de dados, autonomia decisória, modelos, impactos e controles com evidências versionadas. A plataforma não deve exigir exposição pública do dossiê completo, pois pode conter dados pessoais ou segredos comerciais.

\subsection{Quatro perguntas sequenciais}
\begin{enumerate}
\item \textbf{Applicability:} qual jurisdição e quais requisitos são aplicáveis na data e para aquele uso?
\item \textbf{Evidence:} quais registros, avaliações, testes e medidas estão disponíveis?
\item \textbf{Assessment:} há demonstração suficiente de atendimento, descumprimento ou lacuna probatória?
\item \textbf{Change impact:} o que mudou na lei, no produto ou no ambiente desde a última avaliação?
\end{enumerate}

\subsection{Estado de uma obrigação por aplicação}
\begin{tabularx}{\textwidth}{@{}p{.27\textwidth}Y@{}}
\toprule\textbf{Valor} & \textbf{Significado}\\\midrule
\term{Satisfied} & Atendimento demonstrado para o critério e escopo avaliados.\\
\term{NotSatisfied} & Descumprimento demonstrado com base em evidências suficientes.\\
\term{InsufficientEvidence} & Não há prova suficiente para concluir atendimento ou descumprimento.\\
\term{NotApplicable} & Exceção ou ausência de incidência fundamentada.\\
\term{LegalUncertainty} & Aplicação ou interpretação normativa em disputa fundamentada.\\
\term{PendingReview} & Resultado automatizado ainda não validado conforme a política de revisão.\\
\bottomrule
\end{tabularx}

\subsection{OWL, SWRL, SHACL e regras jurídicas}
\begin{tabularx}{\textwidth}{@{}p{.22\textwidth}Y@{}}
\toprule\textbf{Mecanismo} & \textbf{Emprego correto}\\\midrule
OWL 2 & Estrutura conceitual, relações ontológicas e inferências descritivas.\\
SWRL & Regras monotônicas positivas: inferir relações a partir de fatos conhecidos.\\
SHACL & Exigir formato, propriedades mínimas, tipos de dados e integridade dos registros RDF.\\
SPARQL & Consultas ao grafo e detecção declarativa de padrões.\\
LegalRuleML ou motor normativo equivalente & Regras com precedência, exceções, modalidade de obrigação, contextos e vigência.\\
PROV-O & Origem dos fatos e transformações, com trilha auditável.\\
\bottomrule
\end{tabularx}
Padrões relevantes incluem SHACL e PROV-O, publicados pelo W3C, e LegalRuleML, publicado pela OASIS \cite{shacl,prov,legalruleml}.

\subsection{Exemplo didático SWRL: obrigação candidata}
\begin{lstlisting}
hrai:AISystem(?system) ^
hrai:hasRiskClassification(?system, hrai:HighRisk)
  ->
hrai:subjectToRequirement(
  ?system, hrai:RiskManagementRequirement)
\end{lstlisting}
\textbf{Interpretação:} sob um perfil normativo \emph{explicitamente validado}, uma classificação \term{HighRisk} implica uma obrigação de gestão de riscos. Na prática, a aplicabilidade jurídica deve ser previamente filtrada por jurisdição, texto, data, exclusões e papéis. A regra isolada não certifica conformidade.

\subsection{Exemplo SHACL: evidência obrigatória no dossiê}
\begin{lstlisting}
@prefix sh: <http://www.w3.org/ns/shacl#> .
@prefix hrai: <https://example.org/hrai-go#> .

hrai:AssessmentShape a sh:NodeShape ;
  sh:targetClass hrai:ComplianceAssessment ;
  sh:property [
    sh:path hrai:usesLegalSnapshot ;
    sh:minCount 1 ;
  ] ;
  sh:property [
    sh:path hrai:hasAssessmentDate ;
    sh:minCount 1 ;
  ] .
\end{lstlisting}
O SHACL verifica a presença de campos no grafo: não verifica se o \term{LegalSnapshot} escolhido corresponde à interpretação jurídica correta.

\section{Motor DEBATE: avaliação crítica da regulação}
\subsection{A questão epistemológica}
A lei é uma fonte normativa, mas também é objeto de investigação. A plataforma deve perguntar se as escolhas regulatórias são aderentes às necessidades, capacidades e prioridades do país sem converter perspectivas políticas em fatos jurídicos.

\subsection{Taxonomia de tensões}
\begin{longtable}{@{}p{.4\textwidth}p{.52\textwidth}@{}}
\toprule\textbf{Classe canônica} & \textbf{Questão de análise}\\\midrule\endhead
\term{ProtectionInnovationTension} & Proteção de direitos versus custos de experimentação e entrada.\\
\term{SovereigntyOpennessTension} & Autonomia nacional versus circulação internacional de dados, tecnologia e conhecimento.\\
\term{PrivacyUtilityTension} & Proteção de dados versus utilidade social e científica.\\
\term{TransparencySecurityTension} & Transparência versus riscos de segurança e informação sensível.\\
\term{RegulatoryCapacityGap} & Exigências legais versus recursos e competências institucionais.\\
\term{DistributionalImpactTension} & Repartição desigual de ônus, danos e benefícios.\\
\term{RegulatoryCoverageGap} & Risco identificado sem cobertura normativa clara.\\
\term{RegulatoryObsolescence} & Suposta perda de adequação diante de mudanças contextuais ou tecnológicas.\\
\term{JurisdictionalConflict} & Exigências incompatíveis ou tensionadas entre territórios.\\
\term{ImplementationFeasibilityTension} & Viabilidade prática de cumprir e fiscalizar uma política.\\
\bottomrule
\end{longtable}

\subsection{Ciclo de uma controvérsia}
\begin{enumerate}
\item \term{Detect}: identificar uma potencial tensão apoiada em fontes.
\item \term{Substantiate}: exigir indicadores, textos normativos, métodos e evidências contraditórias.
\item \term{Contest}: permitir objeções e contraexemplos.
\item \term{Deliberate}: submeter o problema ao Shadow Council e, quando cabível, a fóruns humanos.
\item \term{Simulate}: comparar alternativas sob hipóteses transparentes.
\item \term{Propose}: publicar hipóteses de melhoria e as condições para testá-las.
\item \term{Observe}: acompanhar desdobramentos empíricos e normativos.
\end{enumerate}

\subsection{Exemplo brasileiro: custo de compliance para pequenas empresas}
\nsource O estudo comparativo sustenta a investigação de tensões entre exigências prévias de alto custo, acesso de pequenas empresas e preservação de direitos fundamentais \cite{comparativo}.\par
\nproposal Um \term{PolicyClaim} questiona se obrigações podem ser graduadas por risco, escala e contexto de uso; outro sustenta que controles básicos contra danos graves não devem variar com porte econômico. Ambos devem apresentar evidência e método, e nenhuma posição se torna ``correta'' apenas por maioria.

\section{Biblioteca comparativa dinâmica por jurisdição}
\subsection{Comparar funções, não rótulos nacionais}
Em vez de concluir genericamente que ``a Europa é rígida'' ou ``os EUA são permissivos'', comparar disposições normativas concretas em contextos datados. Um país pode ter regras muito distintas para IA em recrutamento, crédito, saúde, conteúdo sintético ou proteção de dados.

\begin{longtable}{@{}p{.22\textwidth}p{.29\textwidth}p{.41\textwidth}@{}}
\toprule\textbf{Jurisdição} & \textbf{Porta de entrada} & \textbf{Dimensões comparáveis}\\\midrule\endhead
União Europeia & EU AI Act, GDPR e normas setoriais & Classificações, papéis, obrigações, calendários de aplicação, fiscalização.\\
Brasil & LGPD, normas setoriais, PBIA e propostas legislativas & Setores e competências, proteção de dados, riscos distributivos, capacidade institucional.\\
Reino Unido & Política e autoridades setoriais & Responsabilidade institucional, orientação, fiscalização por setor.\\
Estados Unidos & Instrumentos federais, estaduais e setoriais & Fragmentação territorial, supervisão setorial e deveres aplicáveis.\\
China & Instrumentos relativos a serviços e sistemas específicos & Requisitos por categoria de serviço, segurança e mecanismos de controle.\\
Índia & Política de IA, infraestrutura digital e instrumentos setoriais & Desenvolvimento, inclusão, governança de dados e aplicação prática.\\
\bottomrule
\end{longtable}
\subsection{Unidade de comparação}
Para cada \term{NormativeRequirement}, a biblioteca preservaria \term{jurisdiction}, \term{effectiveInterval}, \term{regulatedActors}, \term{scope}, \term{exceptions}, \term{enforcement}, \term{source} e \term{confidenceOfInterpretation}.\par
A plataforma distinguiria \textbf{equivalência semântica}, \textbf{semelhança funcional} e \textbf{equivalência jurídica}: conceitos parecidos não produzem automaticamente deveres idênticos.

\section{Shadow Council: um conselho deliberativo sintético e plural}
\begin{keybox}
\textbf{Definição.} O \term{ShadowCouncil} é um laboratório de deliberação composto por agentes de IA configurados para \textbf{examinar perspectivas de stakeholders com apoio em posições e evidências documentadas}. Seu objetivo é ampliar a cobertura de argumentos, antecipar questões negligenciadas e testar alternativas. Ele \textbf{não} equivale a consulta pública real, amostragem de opinião, mandato institucional ou votação democrática.
\end{keybox}

\subsection{O que significa ``amostra representativa'' aqui?}
Há duas noções que precisam permanecer separadas:\par
\begin{description}[leftmargin=0pt,style=nextline]
\item[\term{PerspectiveCoverage}] Cobertura de \emph{perspectivas documentadas} relevantes para um tema, por estrato, interesse, papel e contexto. Pode ser estimada e auditada sem pessoas participantes.
\item[\term{PopulationRepresentativeness}] Representatividade estatística de uma população ou de organizações reais. Exige \emph{sampling frame}, procedimento amostral válido, dados e, quando envolver indivíduos, salvaguardas éticas e legais. \textbf{Não é produzida por criar muitos agentes sintéticos.}
\end{description}
\nproposal O MVP terá \term{PerspectiveCoverage} como objetivo explícito. Reivindicações de representatividade populacional só poderão ser feitas em estudos com amostras humanas efetivamente recrutadas e metodologicamente avaliadas.

\subsection{Mapeamento inicial de stakeholders}
O conjunto de estratos varia por política e jurisdição. Exemplo para uma sessão brasileira sobre regulação de IA:\par
\begin{longtable}{@{}p{.25\textwidth}p{.42\textwidth}p{.23\textwidth}@{}}
\toprule\textbf{Estrato} & \textbf{Perspectivas a considerar} & \textbf{Fontes documentais}\\\midrule\endhead
Reguladores e fiscalização & Exequibilidade, competência, inspeção, responsabilização e recursos & Normas, consultas, relatórios e manifestações oficiais\\
Poder público operador & Contratação, serviço público, responsabilidade administrativa e infraestrutura & Planos, auditorias, contratos e documentos técnicos\\
Setor produtivo & Grandes empresas, MPMEs, startups, integradores e desenvolvedores & Consultas, associações e estudos de custo\\
Pesquisa e especialistas & Direito, computação, ciências sociais, economia, ética e políticas públicas & Trabalhos acadêmicos e pareceres públicos\\
Sociedade civil e direitos & Defesa do consumidor, privacidade, pessoas com deficiência, direitos digitais & Consultas públicas e manifestações publicadas\\
Trabalhadores e profissões & Impactos no emprego, autonomia, supervisão e organização do trabalho & Estudos setoriais e posições de entidades\\
Grupos diretamente afetados & Usuários de serviços, beneficiários, candidatos, pacientes e comunidades & Depoimentos autorizados, pesquisas e participação direta\\
Infraestrutura e ecossistema & Energia, computação, dados, cadeia de suprimentos e sustentabilidade & Estudos técnicos, indicadores e dados públicos\\
\bottomrule
\end{longtable}
Estratos não são homogêneos. Por exemplo, empresas pequenas e grandes podem divergir; um órgão regulador e um prestador público podem ter interesses diferentes; e a sociedade civil não se reduz a uma posição única.

\subsection{Estratificação multidimensional}
O conselho deverá considerar pelo menos quatro eixos ao definir cobertura:\par
\begin{itemize}
\item \textbf{Institutional Role:} quem regula, desenvolve, fiscaliza, opera ou é afetado.
\item \textbf{Geographic Coverage:} regiões e diferenças territoriais pertinentes.
\item \textbf{Impact Exposure:} grupos direta ou indiretamente afetados pelo risco estudado.
\item \textbf{Viewpoint Diversity:} posições divergentes dentro do mesmo estrato, não só entre estratos.
\end{itemize}
Outros eixos (porte, setor, acessibilidade, condições culturais, vínculos econômicos) só devem ser adicionados se forem pertinentes e sustentados por fontes. Dados pessoais sensíveis não devem ser inferidos de indivíduos reais para fabricar perfis.

\subsection{Exemplo de composição de uma sessão (não amostra populacional)}
\nexample Para discutir a proporcionalidade da conformidade na regulação brasileira, uma sessão pode usar 20 assentos documentais:\par
\begin{center}
\begin{tabularx}{\textwidth}{@{}Y r p{.29\textwidth}@{}}
\toprule\textbf{Estrato} & \textbf{Assentos} & \textbf{Justificativa de cobertura}\\\midrule
Regulação e fiscalização & 3 & Diferentes competências institucionais\\
Poder público operador & 2 & Órgãos implantadores e compradores\\
Empresas e startups & 3 & Pelo menos um contraste por porte\\
Pesquisa independente & 3 & Disciplinas e métodos distintos\\
Sociedade civil e direitos & 3 & Correntes e temas divergentes\\
Trabalhadores & 2 & Relações de trabalho e profissão\\
Grupos afetados & 2 & Exposição concreta à decisão de IA\\
Infraestrutura e sustentabilidade & 2 & Dependência tecnológica e ambiental\\\midrule
\textbf{Total} & \textbf{20} & \textbf{Quotas ilustrativas, não proporcionais}\\\bottomrule
\end{tabularx}
\end{center}
Não são 20 pessoas. São 20 \term{CouncilSeat} cujas posições podem ser simuladas mediante fontes documentadas. O número não implica peso político, autoridade nem equivalência populacional. Os estratos e quotas devem ser recalculados por tema e auditados quanto à cobertura.

\subsection{Como selecionar e configurar os agentes}
\begin{enumerate}
\item Definir \term{CouncilQuestion} e jurisdição, escopo temporal e impacto esperado.
\item Construir \term{StakeholderMap}: quem é afetado, quem implementa, quem regula e quem pode contestar.
\item Elaborar \term{SamplingFrame}: corpus de posições publicadas, dados de participação, pesquisas e documentos coletados com critérios explícitos.
\item Calcular \term{CoverageMatrix} e lacunas por estrato e perspectiva. Ausência de fonte significa \term{Unrepresented}, não licença para inventar posição.
\item Alocar \term{CouncilSeat} com quota mínima de posições divergentes quando fundamentadas e definir se a sessão é exploratória ou validada.
\item Instanciar \term{PerspectiveAgent} com \term{PerspectiveCharter}: escopo, fontes, limites de inferência, conflitos de interesse e regras de citação.
\item Revisar identidade e autorização para quaisquer agentes que aleguem falar oficialmente em nome de uma instituição.
\end{enumerate}

\begin{warnbox}
\textbf{Proibição metodológica.} Um agente não deve alegar ser uma pessoa real, um regulador investido de poder público, um paciente, um trabalhador ou o representante eleito de um grupo sem autorização. A voz simulada deve aparecer como \textit{synthetic perspective}, com fonte e limitação declaradas.
\end{warnbox}

\subsection{Protocolo deliberativo em seis rodadas}
\begin{figure}[H]
\centering
\begin{tikzpicture}[x=1cm,y=1cm,>=Latex,every node/.style={font=\sffamily\scriptsize},st/.style={draw=blue!65,rounded corners=3pt,fill=light,text width=3.05cm,align=center,minimum height=1.1cm}]
\node[st] (a) at (1.6,3.0) {0. Evidence brief\\and question framing};
\node[st] (b) at (5.5,3.0) {1. Independent\\initial positions};
\node[st] (c) at (9.4,3.0) {2. Cross-examination\\and counterarguments};
\node[st] (d) at (9.4,.7) {3. Structured revision\\and uncertainties};
\node[st] (e) at (5.5,.7) {4. Minority report\\and alternative scenarios};
\node[st] (f) at (1.6,.7) {5. Synthesis, audit\\and human review};
\draw[->,thick,teal] (a)--(b);\draw[->,thick,teal] (b)--(c);\draw[->,thick,teal] (c)--(d);\draw[->,thick,teal] (d)--(e);\draw[->,thick,teal] (e)--(f);
\end{tikzpicture}
\caption{Rodadas propostas para uma sessão do Shadow Council; todas as contribuições permanecem rastreáveis.}
\end{figure}
\textbf{Rodada 0 --- Evidence brief:} apresentar o mesmo dossiê factual a todos, com links, fontes e limites.\par
\textbf{Rodada 1 --- Independent positions:} produzir análise inicial isolada para limitar contaminação e ancoragem entre agentes.\par
\textbf{Rodada 2 --- Cross-examination:} agentes podem contestar interpretações, solicitar provas e indicar premissas omitidas.\par
\textbf{Rodada 3 --- Structured revision:} cada agente pode manter, revisar ou retirar afirmações, com justificativa auditável.\par
\textbf{Rodada 4 --- Minority report:} relatar argumentos minoritários sem dissolvê-los numa média.\par
\textbf{Rodada 5 --- Synthesis and human review:} publicação de síntese, matriz de divergências, evidências e recomendações condicionais.\par
A \term{CouncilSession} deve registrar participantes, versões de modelo, parâmetros, sementes quando disponíveis, documentos, turnos, intervenções do moderador e eventuais conflitos de interesse.

\subsection{Decisão, votação e legitimidade}
\begin{itemize}
\item O council pode agregar \textbf{grau de suporte argumentativo} e \textbf{cobertura}, mas não declarar que 12 agentes sintéticos contra 8 representam a vontade social.
\item Direitos fundamentais não devem ser convertidos em custo compensável numa votação automática.
\item Pesos de votação, caso existam para experimentos, devem ser \textbf{parâmetros de cenário}, publicados e acompanhados de análises de sensibilidade.
\item Parecer oficial, decisões administrativas, sanções e atos legislativos permanecem sob responsabilidade das instituições humanas competentes.
\item Uma consulta pública com pessoas reais constitui \textbf{outro procedimento}, que pode alimentar e validar a plataforma, mas não pode ser substituída por ela.
\end{itemize}

\subsection{Controles de diversidade, dependência e viés de modelos}
Instanciar vinte agentes usando o mesmo modelo e o mesmo corpus pode produzir \textbf{falsa diversidade}. Recomenda-se comparar famílias de modelos quando viável, diversificar corpora legitimamente, isolar as primeiras respostas, aplicar agentes de contestação e medir estabilidade sob variação de prompts e composição. A diversidade dos modelos não garante independência epistemológica: fontes comuns, dados de treinamento compartilhados e arquiteturas parecidas continuam sendo fatores de correlação.

\subsection{Métricas de qualidade do Shadow Council}
\begin{longtable}{@{}p{.33\textwidth}p{.59\textwidth}@{}}
\toprule\textbf{Métrica} & \textbf{O que mede}\\\midrule\endhead
\term{StratumCoverage} & Proporção dos estratos pertinentes com ao menos uma posição documentalmente sustentada.\\
\term{ViewpointCoverage} & Diversidade de teses pertinentes cobertas em cada estrato.\\
\term{EvidenceTraceability} & Fração de alegações verificáveis com referência e recorte de fonte.\\
\term{UnsupportedClaimRate} & Frequência de alegações não sustentadas ou com fonte indevida.\\
\term{DissentRetention} & Preservação de argumentos minoritários relevantes no relatório final.\\
\term{ModelSensitivity} & Alterações de resultado ao mudar prompts, modelos, ordem e composição.\\
\term{ExpertAgreement} & Concordância com avaliação humana independente, por critérios predefinidos.\\
\term{AffectedGroupReview} & Validação ou contestação de conclusões por grupos efetivamente afetados.\\
\term{ConflictOfInterestDisclosure} & Completude de declarações de vínculos e interesses.\\
\bottomrule
\end{longtable}

\section{Formalização do Shadow Council na ontologia}
\subsection{Entidades e relações mínimas}
\begin{lstlisting}
ShadowCouncilSession
  hasQuestion -> CouncilQuestion
  usesSamplingFrame -> SamplingFrame
  hasSeat -> CouncilSeat
  hasRound -> DeliberationRound
  produces -> CouncilReport

CouncilSeat
  allocatedTo -> StakeholderStratum
  assignedAgent -> PerspectiveAgent
  hasSelectionRationale -> EvidenceRecord

PerspectiveAgent
  governedBy -> PerspectiveCharter
  trainedOrConditionedOn -> SourceCollection
  generates -> DeliberationContribution
\end{lstlisting}

\subsection{Exemplo sintético em Turtle}
\begin{lstlisting}
@prefix hrai: <https://example.org/hrai-go#> .
@prefix ex:   <https://example.org/demo#> .

ex:BrazilSMECouncil a hrai:ShadowCouncilSession ;
  hrai:hasQuestion ex:ProportionalComplianceQuestion ;
  hrai:hasSeat ex:SeatSmallEnterprise .

ex:SeatSmallEnterprise a hrai:CouncilSeat ;
  hrai:allocatedTo ex:SMEStakeholderStratum ;
  hrai:assignedAgent ex:SyntheticSMEPerspective .

ex:SyntheticSMEPerspective a hrai:PerspectiveAgent ;
  hrai:representsPerspective ex:SMEStakeholderStratum ;
  hrai:hasSyntheticStatus true ;
  hrai:governedBy ex:SMECharterV1 .
\end{lstlisting}
\nexample Os indivíduos acima não representam uma empresa ou posição reais. Para uma execução, seria necessário conectar \term{SMECharterV1} a documentos públicos identificados e registrar eventuais limitações de cobertura.

\subsection{Regras inferenciais ilustrativas}
Uma regra SWRL pode detectar uma \textit{condição operacional de revisão}, desde que haja um indivíduo de alerta previamente definido. Não é apropriado inferir que uma opinião de um stakeholder é verdadeira por corresponder à posição simulada.
\begin{lstlisting}
hrai:CouncilSeat(?seat) ^
hrai:assignedAgent(?seat, ?agent) ^
hrai:hasSyntheticStatus(?agent, true)
 ->
hrai:requiresSyntheticDisclosure(?seat, true)
\end{lstlisting}
A verificação de que uma \term{CouncilSession} possui assentos documentados e identificados é mais adequadamente implementada por SHACL, pois exige \emph{completude} sob a hipótese de mundo fechado do grafo validado.

\section{Agentes especializados e governança de contribuições}
\begin{longtable}{@{}p{.43\textwidth}p{.49\textwidth}@{}}
\toprule\textbf{Agente} & \textbf{Responsabilidade}\\\midrule\endhead
\term{LegalSourceMonitorAgent} & Monitorar alterações regulatórias e respectivas fontes.\\
\term{NormativeExtractionAgent} & Extrair requisitos candidatos e suas condições.\\
\term{OntologyAlignmentAgent} & Alinhar termos nacionais à ontologia canônica.\\
\term{ComplianceAssessmentAgent} & Avaliar requisitos e evidências de aplicações.\\
\term{NationalContextAgent} & Consolidar indicadores e suas séries temporais.\\
\term{ComparativeRegulationAgent} & Relacionar funções regulatórias entre jurisdições.\\
\term{RegulatoryTensionAgent} & Detectar tensões e gerar questões para investigação.\\
\term{EvidenceReviewAgent} & Contestar afirmações frágeis e verificar sustentação.\\
\term{ShadowCouncilOrchestrator} & Instanciar assentos e organizar rodadas deliberativas.\\
\term{DeliberationAuditAgent} & Avaliar cobertura, rastreabilidade, conflitos e dissenso.\\
\term{PolicySimulationAgent} & Analisar alternativas conforme hipóteses e dados explícitos.\\
\term{ProvenanceAuditAgent} & Preservar linhagem de fontes, regras, modelos e revisões.\\
\bottomrule
\end{longtable}

\subsection{Modelo aberto com níveis de autoridade}
\begin{enumerate}
\item \term{PublicContributor}: pode propor e comentar alegações, mas não alterar fatos validados.
\item \term{VerifiedStakeholder}: identidade institucional verificada; publica posições próprias sob autorização.
\item \term{ExpertReviewer}: pode revisar objetos segundo competência e trilha de auditoria.
\item \term{RegulatorySourceCurator}: valida a correspondência documental com fonte oficial.
\item \term{OntologyMaintainer}: propõe alterações da TBox por processo de governança.
\item \term{SystemOperator}: administra infraestrutura e segurança sem reivindicar autoridade jurídica.
\end{enumerate}
Em todos os casos, proteger contra injeção de prompt em fontes recebidas, adulteração de evidências, agentes que se façam passar por autoridades e exposição de dossiês empresariais ou dados pessoais.

\section{Cenário integrado: análise de uma aplicação e uma tensão brasileira}
\subsection{Pergunta A: o sistema está conforme?}
\nexample Uma organização usa \term{RecruitmentScreeningAI} para ranquear candidatos. O agente \term{VERIFY} identifica dados pessoais e o contexto de recrutamento, consulta o retrato legal em data determinada, produz obrigações candidatas e solicita evidências de revisão, transparência, finalidade e segurança. Onde não há prova suficiente, registra \term{InsufficientEvidence}; não presume ilicitude nem conformidade.

\subsection{Pergunta B: a regulação precisa ser revista?}
O \term{DEBATE} detecta possível \term{ProtectionInnovationTension} em pequenas empresas. O \term{NationalContextAgent} reúne indicadores com período e fonte. O \term{ComparativeRegulationAgent} busca políticas funcionalmente parecidas na UE e no Reino Unido. O \term{ShadowCouncilOrchestrator} cria um conselho com posições documentadas de trabalhadores, empresas, direitos digitais, reguladores, especialistas e candidatos a vagas. O relatório inclui objeções e divergências.

\subsection{Pergunta C: qual alternativa merece teste?}
O \term{FORESIGHT} compara, por exemplo: (i) exigências prévias uniformes; (ii) controles graduados por exposição a risco; (iii) experimentação supervisionada sob condições especiais. As hipóteses devem explicitar custos administrativos, danos potenciais, capacidade de fiscalização e distribuição dos impactos. A comparação não atribui automaticamente uma nota de ``melhor regulação''.

\subsection{Artefatos produzidos no mesmo ciclo}
\begin{itemize}
\item \term{ComplianceAssessment}: avaliação da aplicação, fontes e lacunas probatórias.
\item \term{RegulatoryTensionCase}: questão de política pública e evidências conflitantes.
\item \term{ComparativePolicyReport}: semelhanças e diferenças jurídicas documentadas.
\item \term{ShadowCouncilReport}: argumentos, cobertura e dissenso registrado.
\item \term{PolicyAlternativeAssessment}: cenários e hipóteses testáveis.
\end{itemize}

\section{Esquema de saída canônico (exemplo JSON)}
O objeto abaixo serve apenas para ilustrar a estrutura compartilhada de uma sessão. Não é um \eng{JSON Schema} normativo completo.
\begin{lstlisting}
{
  "sessionId": "council-br-001",
  "jurisdiction": "BR",
  "mode": "synthetic-perspective-deliberation",
  "question": "Which proportionality criteria should be tested?",
  "legalSnapshotId": "legal-snapshot-br-2026-11-11",
  "samplingFrame": {
    "type": "documentary",
    "populationRepresentative": false,
    "selectionMethod": "purposive-stratified-coverage",
    "coverageLimitations": ["positions are not public opinion"]
  },
  "seats": [
    {
      "seatId": "sme-seat-01",
      "stratum": "small-enterprises",
      "agentId": "synthetic-sme-01",
      "synthetic": true,
      "sourceIds": ["verified-source-001"]
    }
  ],
  "arguments": [],
  "counterArguments": [],
  "minorityReport": [],
  "evidenceReviewStatus": "pending",
  "institutionalEndorsement": false,
  "humanReviewStatus": "pending"
}
\end{lstlisting}

\section{Validação científica e critérios de qualidade}
\subsection{Questões de pesquisa}
\begin{tabularx}{\textwidth}{@{}p{.13\textwidth}Y@{}}
\toprule\textbf{RQ} & \textbf{Pergunta}\\\midrule
RQ1 & A atualização contínua melhora a precisão e a tempestividade das avaliações de conformidade?\\
RQ2 & A estrutura ontológica facilita comparações regulatórias reproduzíveis entre países?\\
RQ3 & O Shadow Council amplia a diversidade de argumentos relevantes em comparação com uma análise de agente único?\\
RQ4 & A estratificação documental aumenta a cobertura sem produzir falsa alegação de representatividade?\\
RQ5 & Deliberação estruturada e auditoria de proveniência reduzem alegações sem suporte e preservam dissensos?\\
RQ6 & O contexto socioeconômico altera materialmente a escolha entre alternativas regulatórias?\\
\bottomrule
\end{tabularx}

\subsection{Desenho experimental inicial}
Comparar três condições em casos brasileiros previamente selecionados: \textbf{(A)} análise por agente único; \textbf{(B)} múltiplos agentes sem estratificação; \textbf{(C)} Shadow Council estratificado, com fontes verificadas, debate e auditoria. Uma banca humana multidisciplinar, independente dos resultados automatizados, define rubricas e examina argumentos. Registrar cobertura de teses, precisão de fontes, falsos positivos de lacunas, retenção de dissenso e custo. Repetir com diferentes composições, prompts e famílias de modelos.

\subsection{Critérios mínimos para aceite}
\begin{itemize}
\item Todos os resultados jurídicos publicados possuem dispositivo, versão e status normativo vinculados.
\item Nenhuma proposição pendente é apresentada como obrigação legal vigente.
\item Toda interpretação contestável pode ser contestada sem apagamento do registro anterior.
\item Toda perspectiva sintética é identificada como tal, com escopo e documentação.
\item Resultados do Shadow Council incluem cobertura, limitações, posições minoritárias e revisão humana.
\item Avaliações de compliance são reproduzíveis a partir de versões das regras, entradas e evidências.
\end{itemize}

\section{Roadmap incremental de implementação}
\begin{longtable}{@{}p{.2\textwidth}p{.47\textwidth}p{.25\textwidth}@{}}
\toprule\textbf{Fase} & \textbf{Entregas} & \textbf{Critério de saída}\\\midrule\endhead
MVP 1 --- NORM & Ontologia modular, fontes oficiais brasileiras, histórico legislativo, versões e proveniência. & Mudanças rastreáveis e reprocessamento determinístico.\\
MVP 2 --- VERIFY & Dossiê de aplicação, regras de incidência, SHACL e avaliações versionadas. & Casos avaliados com revisão jurídica e trilha probatória.\\
MVP 3 --- DEBATE & Tensões, claims, contra-argumentos, atores externos e biblioteca BR/UE/UK. & Relatórios comparativos contestáveis e auditáveis.\\
MVP 4 --- SHADOW COUNCIL & Estratos, frames documentais, assentos sintéticos, rodadas e relatórios de minoria. & Melhor cobertura verificada contra baseline; nenhum falso mandato de representação.\\
MVP 5 --- FORESIGHT & Alternativas, cenários, indicadores e acompanhamento posterior. & Hipóteses reprodutíveis e incerteza publicada.\\
\bottomrule
\end{longtable}

\subsection{Tecnologia: decisões reversíveis}
Uma implementação possível inclui repositório RDF ou grafo semântico, buscas textuais e vetoriais sobre textos normativos, filas de eventos, mecanismos de inferência e APIs. A escolha entre triplestore, banco relacional com suporte RDF, grafo de propriedades ou combinação deve ser determinada por testes de escala, versionamento, SPARQL, proveniência, licença e custo operacional. A ontologia é estável conceitualmente; o fornecedor de banco não precisa ser.

\section{Riscos e salvaguardas de governança}
\begin{longtable}{@{}p{.31\textwidth}p{.59\textwidth}@{}}
\toprule\textbf{Risco} & \textbf{Salvaguarda proposta}\\\midrule\endhead
Alucinação regulatória & Fonte oficial obrigatória, status temporal, revisão por criticidade e linguagem de incerteza.\\
Representatividade artificial & Distinguir cobertura documental de amostragem real; impedir simulação de identidade não autorizada.\\
Captura por stakeholders & Declaração de financiamento e conflito, transparência de critérios e preservação do contraditório.\\
Domínio de um modelo & Teste de sensibilidade a modelos e prompts, críticas adversariais e comparação com humanos.\\
Divulgação de dados sensíveis & Isolamento multitenant, controle de acesso e publicação agregada conforme necessidade.\\
Automação de decisão política & Conselho sintético é ferramenta consultiva; decisões formais são humanas e institucionais.\\
Obsolescência das regras & Monitoramento da fonte oficial, eventos, snapshots e recomputação das avaliações afetadas.\\
Simulação causal infundada & Hipóteses explícitas, limites de validade, intervalos e avaliação empírica antes de inferir efeitos.\\
\bottomrule
\end{longtable}

\section{Síntese e posicionamento}
O HRAI-GO pode ser conceituado como uma \textbf{plataforma aberta de inteligência regulatória baseada em conhecimento}, com agentes especializados e participação estruturada dos stakeholders. A ontologia organiza fatos, requisitos e controvérsias; os mecanismos normativos examinam aplicações; a biblioteca internacional contextualiza alternativas; e o Shadow Council amplia a cobertura argumentativa para preparar deliberações humanas informadas.

Seu valor singular não está em automatizar o debate democrático ou declarar a ``melhor'' política. Está em produzir \textbf{rastreabilidade, comparabilidade, capacidade de contestação, memória institucional e hipóteses públicas verificáveis} em um processo contínuo. A distinção entre regulação vigente, evidências observadas e sugestões normativas é a condição fundamental para essa proposta fazer sentido.

\begin{notebox}
\textbf{Princípio de projeto.} \emph{Evidence before inference; inference before recommendation; recommendation before human decision.} Em português: primeiro a evidência, depois a inferência, depois a recomendação e, por fim, a decisão humana competente.
\end{notebox}

\appendix
\section{Glossário bilíngue seletivo}
\begin{longtable}{@{}p{.4\textwidth}p{.52\textwidth}@{}}
\toprule\textbf{Termo canônico} & \textbf{Explicação}\\\midrule\endhead
\term{RegulatoryIntelligence} & Inteligência sobre normas, instituições, mudanças e efeitos regulatórios.\\
\term{LegalSnapshot} & Recorte datado e versionado do quadro jurídico conhecido.\\
\term{NormativeRequirement} & Requisito jurídico interpretado e vinculado à origem.\\
\term{ComplianceAssessment} & Avaliação de incidência normativa e evidências de uma aplicação.\\
\term{RegulatoryGap} & Lacuna potencial que necessita confirmação e fundamentação.\\
\term{StakeholderStratum} & Conjunto de atores reunidos por critério explícito de participação.\\
\term{SamplingFrame} & Base de referência para selecionar posições ou participantes.\\
\term{ShadowCouncil} & Deliberação simulada por perspectivas sintéticas documentadas.\\
\term{PerspectiveCoverage} & Cobertura verificável de perspectivas pertinentes ao problema.\\
\term{DissentRetention} & Preservação de argumentos minoritários na síntese final.\\
\term{Provenance} & Histórico de origem, processamento e autoria de uma afirmação ou dado.\\
\term{BitemporalModel} & Representação separada de tempo de validade e tempo de registro.\\
\bottomrule
\end{longtable}

\section{Questões de competência para testar a ontologia}
\begin{enumerate}
\item Quais obrigações são aplicáveis ao \term{AISystem X} no Brasil em uma data indicada?
\item Que dispositivos mudaram entre dois \term{LegalSnapshot}?
\item Que aplicações devem ser reavaliadas após uma alteração normativa?
\item Que requisitos funcionalmente semelhantes existem em outras jurisdições?
\item Qual a origem de cada conclusão de \term{ComplianceAssessment}?
\item Quais evidências contrariam a hipótese de uma \term{RegulatoryGap}?
\item Quais segmentos de stakeholders não foram contemplados no \term{SamplingFrame}?
\item Um \term{ShadowCouncilReport} preservou argumentos minoritários fundamentados?
\item Quanto os resultados mudam ao alterar os critérios de composição do council?
\item Quais alternativas regulatórias permanecem inconclusivas por falta de dados?
\end{enumerate}

\section{Referências e rastreabilidade}
\begin{thebibliography}{20}
\bibitem{guia} \textit{Guia de Conformidade e Governança para Sistemas de Inteligência Artificial de Alto Risco (HRAI)}. Estudo fornecido pelo solicitante, 23 p. Fonte de temas e conceitos de compliance; não usado como prova de vigência de normas.
\bibitem{comparativo} \textit{Análise Comparativa e Diretrizes Estratégicas para a Governança e Plano de Inteligência Artificial no Brasil: Alinhamento com o AI Act, Modelos Internacionais e Especificidades Nacionais}. Estudo fornecido pelo solicitante, 16 p. Fonte de hipóteses estratégicas para investigação.
\bibitem{brainstorm} \textit{Brainstorming: tipos de artefatos técnicos que não fiquem datados}. Material fornecido pelo solicitante, 2 p. Fonte de propostas de artefatos permanentes.
\bibitem{camara} Câmara dos Deputados. \textit{PL 2338/2023: ficha de tramitação}. \url{https://www.camara.leg.br/proposicoesWeb/fichadetramitacao?idProposicao=2487262}. Consultado em 11 nov. 2026.
\bibitem{eutimeline} European Commission. \textit{Timeline for the Implementation of the EU AI Act}. \url{https://ai-act-service-desk.ec.europa.eu/en/ai-act/eu-ai-act-implementation-timeline}. Consultado em 11 nov. 2026.
\bibitem{shacl} W3C. \textit{Shapes Constraint Language (SHACL)}. \url{https://www.w3.org/TR/shacl/}. Recomendação técnica.
\bibitem{prov} W3C. \textit{PROV-O: The PROV Ontology}. \url{https://www.w3.org/TR/prov-o/}. Recomendação técnica.
\bibitem{legalruleml} OASIS. \textit{LegalRuleML Core Specification Version 1.0}. \url{https://docs.oasis-open.org/legalruleml/legalruleml-core-spec/v1.0/}. Especificação técnica.
\bibitem{oecd} OECD. \textit{OECD AI Policy Observatory}. \url{https://oecd.ai/}. Fonte para mapeamento de políticas e indicadores.
\bibitem{unesco} UNESCO. \textit{AI Readiness Assessment Methodology}. \url{https://www.unesco.org/ethics-ai/en/ram}. Referência para avaliação nacional multidimensional.
\end{thebibliography}
\end{document}
